\chapter{Classical Numerical Methods}
\label{ch:numerics}

\epigraph{\itshape A good approximation is one that holds together for a
while, and then breaks in a way you understand.}{Folklore in numerical
analysis}

A SciML practitioner cannot escape numerical analysis: it furnishes
the training data, the baselines, and the inductive biases for many
neural-network methods.

\section{Time integration of ODEs}
Given $\dot{\bm{x}}=\bm{f}(\bm{x},t)$, denote the time step $h$.

\subsection{Explicit Runge--Kutta}
The classical $s$-stage RK scheme is
\begin{equation}
\bm{k}_i = \bm{f}\!\left(\bm{x}_n + h\sum_{j=1}^{i-1}a_{ij}\bm{k}_j,\; t_n+c_ih\right),\quad
\bm{x}_{n+1}=\bm{x}_n+h\sum_{i=1}^s b_i\bm{k}_i.
\end{equation}
RK4 ($s=4$) achieves order four with the well-known Butcher tableau.
Adaptive Dormand--Prince (DP5) embeds an order-4 estimator inside an
order-5 step for error control.

\subsection{Implicit and stiff solvers}
Stiff systems --- where the spectrum of the Jacobian spans many orders
of magnitude --- demand implicit methods. Backward differentiation
formulas (BDF), implicit RK (Radau, SDIRK), and Rosenbrock methods are
standard. The cost is a (possibly nonlinear) solve at each step.

\begin{insight}{Why ``stiff'' bites neural ODEs}
Neural ODEs use \emph{black-box} adaptive solvers and back-propagate
through them. When the learned dynamics become stiff, solver step
counts (and wall time) explode. Mitigations include semi-implicit
adjoints, regularising stiffness, or restricting parameterisations
(\cref{ch:nodes}).
\end{insight}

\subsection{Symplectic integrators}
For Hamiltonian systems, preserving $\omega$ matters more than local
order. \emph{St\"ormer--Verlet} and higher-order \emph{Yoshida}
splittings have bounded energy error over exponentially long times --
critical for n-body astrophysics.

\begin{recipe}{Choosing a time integrator}
\begin{enumerate}
\item Non-stiff, non-conservative: RK4, DP5, Tsit5.
\item Stiff: BDF, Radau IIA, Rosenbrock--Wanner.
\item Hamiltonian / long-horizon energy fidelity: Verlet, Yoshida-4.
\item Stochastic: Euler--Maruyama (order $\tfrac12$), Milstein
(order $1$), stochastic Runge--Kutta.
\end{enumerate}
\end{recipe}

\section{Spatial discretisation of PDEs}
\subsection{Finite differences (FD)}
Replace derivatives by stencils on a uniform grid:
$\partial_x u_i \approx (u_{i+1}-u_{i-1})/(2h)$. Easy to implement,
limited geometric flexibility, suffer at sharp gradients without
upwinding or limiters.

\subsection{Finite volumes (FV)}
Integrate conservation laws over control volumes; fluxes at faces are
reconstructed via Godunov / Roe / HLLC solvers. Native handling of
shocks and discrete conservation. Backbone of CFD codes.

\subsection{Finite elements (FE)}
Project the weak form $\int_\Omega(\Lop u)v\,\dd\Omega=\int_\Omega
fv\,\dd\Omega$ onto a finite-dimensional basis $\{\phi_j\}$.
Galerkin / Petrov--Galerkin / DG variants. Excellent on complex
geometries.

\subsection{Spectral methods}
Expand $u(\bm{x})=\sum_k \hat u_k\phi_k(\bm{x})$ in a global basis
(Fourier on tori, Chebyshev on intervals). Exponential convergence
for analytic solutions; foundational for Fourier neural operators
(\cref{sec:fno}).

\subsection{Mesh-free / particle methods}
Smoothed particle hydrodynamics (SPH), radial basis functions (RBF),
and reproducing kernel particle methods. Neural networks themselves can
be viewed as a mesh-free basis.

\section{Linear algebra and Krylov methods}
Most PDE discretisations reduce to large sparse linear systems
$\bm{A}\bm{u}=\bm{b}$. Conjugate gradient (CG), GMRES, BiCGStab, and
multigrid are workhorses. Preconditioning ($\bm{M}^{-1}\bm{A}\bm{u}=
\bm{M}^{-1}\bm{b}$) is often the difference between a usable and
unusable method.

\section{Adjoint and sensitivity analysis}
For gradient-based learning of physical parameters, the
\emph{continuous adjoint} computes
$\dd L/\dd\mu$ at cost independent of $\dim\mu$ via a backward ODE:
\begin{equation}
\dot{\bm{\lambda}} = -\bm{\lambda}^\T \frac{\partial \bm{f}}{\partial \bm{x}} - \frac{\partial \ell}{\partial \bm{x}},\qquad
\frac{\dd L}{\dd \mu} = \int_0^T \bm{\lambda}^\T \frac{\partial \bm{f}}{\partial \mu}\,\dd t.
\end{equation}
The discrete adjoint differentiates the solver itself; both are special
cases of automatic differentiation (Appendix~A).

\section{Stability, consistency, convergence}
The Lax equivalence theorem ties these together for linear problems:
\emph{a consistent finite-difference scheme is convergent if and only
if it is stable}. The CFL condition $|c|h_t\le h_x$ for advection,
the stiffness-driven step restriction for explicit parabolic
schemes, and the von Neumann amplification analysis are all canonical.
These constraints transfer almost verbatim to neural-network surrogates
that mimic explicit time-stepping.

\begin{pitfall}{Surrogate ``solvers'' that violate CFL}
A learned time-stepper trained on $\Delta t$ slightly larger than the
finest physical time scale may appear accurate in-distribution but
diverge catastrophically when rolled out --- exactly the failure mode
of an unstable explicit scheme. \emph{Train-time rollouts} or
implicit-style architectures help.
\end{pitfall}

\section{Reduced-order models (ROMs)}
For high-dimensional systems, project onto a low-dimensional subspace.
Proper Orthogonal Decomposition (POD) builds modes from snapshots via
SVD. Galerkin or Petrov--Galerkin projection yields a small ODE on the
modes. ROMs are often the bridge between full-order numerics and
neural-operator surrogates --- and provide Koopman/DMD with their data
(\cref{ch:operator-theoretic}).

\section{Summary}
Classical numerics gives us a vocabulary --- consistency, stability,
order, conservation, symplecticity --- that survives intact in modern
SciML, even when the basis functions are replaced by neural
networks. The interesting question is not whether to abandon classical
methods, but \emph{which of their structural properties to preserve}
inside a learned model.
